\chapter{Concepts}

\section*{Continuous Functions}

\subsection*{Forward Difference}
According to Tyler's theorem,\\[0.25em]

\vspace{-1.5\baselineskip}
$\begin{aligned}
   & f(x_{i+1}) = f(x_i) + \red{f^{'}(x_i)}(\Delta x) + \frac{f^{''}(x_i)}{2!}(\Delta x)^2+\ldots \quad \quad [\text{where } x = x_{i+1} - x_i] \\
  \Rightarrow & \red{f^{'}(x_i)} = \frac{f(x_{i+1}) - f(x_i)}{\Delta x} - \frac{f^{''}(x_i)}{2!}(\Delta x) +\ldots \\
  \Rightarrow & f^{'}(x_i) = \frac{f(x_{i+1}) - f(x_i)}{\Delta x} + {0}(\Delta x) \quad \quad [\text{Neglecting higher-order terms}]
\end{aligned}$

\subsection*{Backward Difference}
According to Tyler's theorem,\\[0.25em]

\vspace{-1.5\baselineskip}
$\begin{aligned}
   & f(x_{i-1}) = f(x_i) - \red{f^{'}(x_i)}(\Delta x) + \frac{f^{''}(x_i)}{2!}(\Delta x)^2+\ldots \quad \quad [\text{where } x = x_i - x_{i-1}] \\
   \Rightarrow & \red{f^{'}(x_i)} = \frac{f(x_i) - f(x_{i-1})}{\Delta x} + \frac{f^{''}(x_i)}{2!}(\Delta x) +\ldots \\
   \Rightarrow & f^{'}(x_i) = \frac{f(x_i) - f(x_{i-1})}{\Delta x} + {0}(\Delta x) \quad \quad [\text{Neglecting higher-order terms}]
\end{aligned}$

\subsection*{Central Difference}
According to Tyler's theorem,\\[0.25em]

\vspace{-1.5\baselineskip}
$\begin{aligned}
   & f(x_{i+1}) = f(x_i) + \red{f^{'}(x_i)}(\Delta x) + \frac{f^{''}(x_i)}{2!}(\Delta x)^2 + \frac{f^{'''}(x_i)}{3!}(\Delta x)^3 + \ldots \quad \quad [\text{where } x = x_{i+1} - x_i]
   \quad \quad \quad \quad \quad \quad (1)\\
   & f(x_{i-1}) = f(x_i) - \red{f^{'}(x_i)}(\Delta x) + \frac{f^{''}(x_i)}{2!}(\Delta x)^2 - \frac{f^{'''}(x_i)}{3!}(\Delta x)^3 + \ldots \quad \quad [\text{where } x = x_i - x_{i-1}]
   \quad \quad \quad \quad \quad \quad (2)\\[1em]
\end{aligned}$\\
% 
Subtracting equation $(2)$ from equation $(1)$\\[0.5em]
$\begin{aligned}
   \Rightarrow & f(x_{i+1}) - f(x_{i-1}) = 2\red{f^{'}(x_i)}(\Delta x) + \frac{2f^{'''}(x_i)}{3!}(\Delta x)^3 +\ldots\\
   \Rightarrow & \frac{f(x_{i+1}) - f(x_{i-1})}{2(\Delta x)} = \red{f^{'}(x_i)} + \frac{f^{'''}(x_i)}{3!}(\Delta x)^2+\ldots\\
   \Rightarrow & \red{f^{'}(x_i)} = \frac{f(x_{i+1}) - f(x_{i-1})}{2(\Delta x)} - \frac{f^{'''}(x_i)}{3!}(\Delta x)^2+\ldots\\
   \Rightarrow & f^{'}(x_i) = \frac{f(x_{i+1}) - f(x_{i-1})}{2(\Delta x)} - {0}(\Delta x)^2 \quad \quad [\text{Neglecting higher-order terms}]\\
\end{aligned}$

\pagebreak
\section*{Newton-Raphson}
\begin{enumerate}[label=\textbf{\textit{Step-\arabic*}} :, left=0pt, itemindent=0pt, parsep=0.5em]
   \item Calculate  $f^{'}(x)$ symbolically.
   \item Choose an initial guess of the root \red{$x_i$}, to estimate the new
   value of the root \red{$x_{i+1}$} as,\\[0.5em]
   $x_{i+1} = x_i - \frac{f(x_i)}{f^{'}(x_i)}$
   \vspace{0.2em}
   \item Find the absolute relative approximate error |$\in_{a}$| as \\[0.5em]
   $|\in_{a}| = \left| \frac{x_{i+1} - x_i}{x_{i+1}} \right| \times 100$
   \vspace{0.2em}
   \item Compare the absolute relative approximate error with the pre-specified relative error tolerance |$\in_{s}$|.
   \item If |$\in_{a}$| $\geq$ |$\in_{s}$|, then Go to Step 2 using new
   estimate of the root.\\
   Else stop the algorithm.
\end{enumerate}

\section*{Secant Method}
\begin{enumerate}[label=\textbf{\textit{Step-\arabic*}} :, left=0pt, itemindent=0pt, parsep=0.5em]
   \item Calculate the next estimate of the root from two initial guesses \red{$x_i$} and \red{$x_{i-1}$},\\[0.5em]
   $x_{i+1} = x_i - \frac{f(x_i)(x_i - x_{i-1})}{f(x_i) - f(x_{i-1})}$
   \vspace{0.2em}
   \item Find the absolute relative approximate error |$\in_{a}$| as \\[0.5em]
   $|\in_{a}| = \left| \frac{x_{i+1} - x_i}{x_{i+1}} \right| \times 100$
   \vspace{0.2em}
   \item Compare the absolute relative approximate error with the pre-specified relative error tolerance |$\in_{s}$|.
   \item If |$\in_{a}$| $\geq$ |$\in_{s}$|, then Go to Step 1 using new
   estimate of the root.\\
   Else stop the algorithm.
\end{enumerate}

\section*{Gaussian Elimination}
\textbf{\textit{Forward Elimination}} :
\vspace{-0.5\baselineskip}
\begin{itemize}[parsep=0.25em]
   \item Transform the system into an upper triangular matrix using elementary row operations.
   \item To simplify the system so each equation has fewer variables/unknowns than the one above.
\end{itemize}

\textbf{\textit{Back Substitution}} :
\vspace{-0.5\baselineskip}
\begin{itemize}[parsep=0.25em]
   \item Solve the simplified upper triangular system from the bottom row upward.
   \item To find the values of the variables/unknowns in the original system of equations.
\end{itemize}


\section*{LU Decomposition}
Given the system of equations $[A][X] = [C]$, where:
\vspace{-0.5\baselineskip}
\begin{itemize}[parsep=0.25em]
   \item $[A]$ is the coefficient matrix; \quad
   \gray{$\bullet$} \ $[X]$ is the unknown matrix; \quad
   \gray{$\bullet$} \ $[C]$ is the constant matrix; \quad
\end{itemize}
   
Decompose $[A]$ into lower and upper triangular matrices: Such that $[A] = [L][U]$. 
Where:
\vspace{-0.5\baselineskip}
\begin{itemize}[parsep=0.25em]
   \item $[U]$ : upper triangular matrix; \ obtained by applying elementary row operations on $[A]$. (forward elimination)
   \item $[L]$ : lower triangular matrix; \ obtained by the multipliers used in the forward elimination steps.
\end{itemize}


\textbf{\textit{Forward substitution}}: Solve $[L][Z] = [C]$ for $[Z]$ \\[0.5em]
\textbf{\textit{Back substitution}}: Solve $[U][X] = [Z]$ for $[X]$
   
The solution $[X]$ gives the values of the unknowns in the original system.


% \textbf{\textit{Forward Elimination}} :
% \vspace{-0.5\baselineskip}
% \begin{itemize}[parsep=0.25em]
%    \item Transform the system into an \orange{upper triangular matrix} using elementary row operations. $(A \to U)$
%    \item Parallelly compute the \orange{lower triangular matrix} using the multipliers used in the forward elimination steps. $(L)$
% \end{itemize}

% \textbf{\textit{Back Substitution}} : \\[0.2em]
% Back substitution is a method used to solve a system of linear equations where the coefficient matrix is upper triangular.\\
% To solve the system of equations, we start from the last equation and substitute the known values into the previous equations.


\section*{Trapezoidal Rule}

\textbf{\textit{Single Segment}} :\\[0.5em]
$\begin{aligned}
   \int_{a}^{b} f(x) dx & \approx \text{Area of trapezoid} \\[0.5em]
   & = \ \frac{1}{2}(\text{Sum of parallel sides})(\text{height}) \\[0.5em]
   & = \ \frac{1}{2}\red{\left[\right.}f(b) + f(a)\red{\left.\right]}(b-a) \\[0.5em]
   & = \ \frac{(b-a)}{2}(f(a) + f(b))
\end{aligned}$


\textbf{\textit{Multiple Segment}} :\\[0.5em]
$\begin{aligned}
   \int_{a}^{b} f(x) dx \ &\approx \ \frac{h}{2}\red{\left[\right.}f(x_0) + 2f(x_1) + 2f(x_2) + 2f(x_{n-1}) + f(x_n)\red{\left.\right]} \qquad \qquad \text{where, } h=\frac{b-a}{n} \\[0.5em]
   & = \ \frac{h}{2}\left[f(x_0) + 2\left\{\sum_{i=1}^{n-1}f(x_i)\right\} + f(x_n)\right]\\
   & = \ \frac{b-a}{2n}\left[f(a) + 2\left\{\sum_{i=1}^{n-1}f(a+ih)\right\} + f(b)\right]
\end{aligned}$


\section*{Simpson's 1/3$^\text{rd}$ Rule}

\textbf{\textit{Multiple Segment}} :\\[0.5em]
$\begin{aligned}
   \int_{a}^{b} f(x) dx \ &= \ \frac{h}{3}\red{\left[\right.}f(x_0) + 4\left[f(x_1) + f(x_3) + \ldots + f(x_{n-1})\right] + 2\left[f(x_2) + f(x_4) + \ldots + f(x_{n-2})\right] + f(x_n)\red{\left.\right]} \\
   & = \ \frac{h}{3}\left[f(x_0) + 4\left\{\sum_{\substack{i = 1 \\ i \textit{ odd}}}^{n-1}f(x_i)\right\} + 2\left\{\sum_{\substack{i = 2 \\ i \textit{ even}}}^{n-2}f(x_i)\right\} + f(x_n)\right]\\
   & = \ \frac{b-a}{3n}\left[f(a) + 4\left\{\sum_{\substack{i = 1 \\ i \textit{ odd}}}^{n-1}f(a+ih)\right\} + 2\left\{\sum_{\substack{i = 2 \\ i \textit{ even}}}^{n-2}f(a+ih)\right\} + f(b)\right]
\end{aligned}$

\section*{Newton Divided Difference}
\textbf{\textit{Quadratic}} :\\[0.5em]
Given $(x_0,y_0)$, $(x_1, y_1)$, and $(x_2,y_2)$, fit a quadratic interpolant through the data. $f_2(x) = b_0 + b_1 (x-x_0) + b_2 (x-x_0)(x-x_1)$ \\[0.5em]
$\begin{aligned}
   & b_0 = f(x_0) \\
   & b_1 = \frac{f(x_1) - f(x_0)}{x_1 - x_0} \\[0.5em]
   & b_2 = \frac{\frac{f(x_2) - f(x_1)}{x_2 - x_1} - \frac{f(x_1) - f(x_0)}{x_1 - x_0}}{x_2 - x_0}
\end{aligned}$

\pagebreak
\section*{Euler's Method}
Use Euler's method with h=0.1 to find approximate values for the solution of the initial value problem\\[0.2em]
$\frac{dy}{dx}+2y= x^3e^{-2x}, y(\red{0})=\blue{1}$
at x=0.1,0.2,0.3



\textbf{Solution:}
\begin{multicols}{2}
\noindent
\raggedcolumns
Given,\\[0.5em]
$\frac{dy}{dx}+2y= x^3e^{-2x}, y(\red{0})=\blue{1}$

Rewriting the equation in differential form,\\[0.5em]
$\frac{dy}{dx} = x^3e^{-2x} - 2y$

From given equation, $h = 0.1$ \quad $\red{x_0 = 0}, \ \blue{y_0 = 1}$


As we know, \ $x_{i+1}=x_i+h$\\
Now, \
$\begin{aligned}
   & x_1=x_0+h=0+0.1=0.1\\
   & x_2=x_1+h=0.1+0.1=0.2\\
   & x_3=x_2+h=0.2+0.1=0.3\\
\end{aligned}$


\columnbreak

$\begin{aligned}
   y_1 & = y_0 + f(x_0,y_0) \times h \\
   & = 1 + f(\red{0},\blue{1}) \times 0.1 \\
   & = 0.8
\end{aligned}$

$\begin{aligned}
   y_2 & = y_1 + f(x_1,y_1) \times h \\
   & = 0.8 + f(\red{0.1},\blue{0.8}) \times 0.1 \\
   & = 0.640081
\end{aligned}$

$\begin{aligned}
   y_3 & = y_2 + f(x_2,y_2) \times h \\
   & = 0.640081 + f(\red{0.2},\blue{0.640081}) \times 0.1 \\
   & = 0.5126017
\end{aligned}$
\end{multicols}



\section*{Gaussian Elimination}
\textbf{\textit{Question:}} 
Solve the system of linear equations using Naive Gaussian Elimination,\\[0.5em]
$\begin{aligned}
   & 25x^3 + 5x^2 + x = 106.8 \\
   & 64x^3 + 8x^2 + x = 177.2 \\
   & 144x^3 + 12x^2 + x = 279.2 \\
\end{aligned}$


\textbf{\textit{Solution}}:\\[0.5em]
Writing the system of linear equations in $[A][x]=[C]$ form,

\begin{multicols}{2}
\noindent
\raggedcolumns

\textit{Now},\\[0.5em]
$\begin{aligned}
&
\left(
\begin{array}{ccc|c}
25 & 5 & 1 & 106.8 \\
64 & 8 & 1 & 177.2 \\
144 & 12 & 1 & 279.2
\end{array}
\right)
\begin{array}{l}
\end{array}\\ 
\Rightarrow &
\left(
\begin{array}{ccc|c}
25 & 5 & 1 & 106.8 \\
0 & -4.8 & -1.56 & -96.208 \\
0 & -16.8 & -4.76 & -335.968
\end{array}
\right)
\begin{array}{l}
\\
R_2=R_2 - \frac{64}{25}R_1\\[0.5em]
R_3=R_3 - \frac{144}{25}R_1\\
\end{array}\\
\Rightarrow &
\left(
\begin{array}{ccc|c}
25 & 5 & 1 & 106.8 \\
0 & -4.8 & -1.56 & -96.208 \\
0 & 0 & 0.7 & 0.76
\end{array}
\right)
\begin{array}{l}
\\\\
R_3=R_3 - \frac{-16.8}{-4.8}R_2\\
\end{array}\\
\Rightarrow &
\left(
\begin{array}{ccc|c}
1 & 0.2 & 0.04 & 4.272 \\
0 & 1 & 0.325 & 20.043 \\
0 & 0 & 1 & 1.0857
\end{array}
\right)
\begin{array}{l}
R_1=R_1 \div 25\\
R_2=R_2 \div -4.8\\
R_3=R_3 \div 0.7
\end{array}
\end{aligned}$

\columnbreak

Expressing this in equations form,\\[0.2em]
$\begin{aligned}
   & x_1 + 0.2x_2 + 0.04x_3 = 4.272 \\
   & x_2 + 0.325x_3 = 20.043 \\
   & x_3 = 1.0857
\end{aligned}$

\vspace{1em}

By back substitution,\\[0.2em]
$\begin{aligned}
   \therefore x_3 &= 1.0857
   \\[1em]
   % 
   \therefore x_2 &= 20.043 - 0.325x_3\\
   &= 20.043 - 0.325(1.0857) = 19.6901
   \\[1em]
   % 
   \therefore x_1 &= 4.272 - 0.04x_3 - 0.2x_2\\
   &= 4.272 - 0.04(1.0857) - 0.2(19.6901) = 0.2906
\end{aligned}$
\end{multicols}


\pagebreak
\section*{Newton-Raphson}
\textbf{\textit{Question:}}
Use the Newton-Raphson method to find the root of the equation $x^3 - 0.165x^2 + 3.993 \times 10^{-4} = 0$.\\
Use 0.05 as the initial guess.

\begin{multicols}{2}
\noindent
\raggedcolumns
\textbf{\textit{Solution:}}\\[0.5em]
Let $f(x) = x^3 - 0.165x^2 + 3.993 \times 10^{-4}$\\[0.5em]
Then, \ $f^{'}(x) = 3x^2 - 0.33x$\\[0.5em]
Let, $x_0 = 0.05$

\textit{Iteration-1}:\\[0.5em]
$\begin{aligned}
   x_{1} &= x_0 - \frac{f(x_0)}{f^{'}(x_0)} \\
   &= 0.05 - \frac{f(0.05)}{f^{'}(0.05)} \\
   &= 0.05 - \frac{(0.05)^3 - 0.165(0.05)^2 + 3.993 \times 10^{-4}}{3(0.05)^2 - 0.33(0.05)} \\
   &= 0.06242
\end{aligned}$

\columnbreak
\textit{Iteration-2}:\\[0.5em]
$\begin{aligned}
   x_{2} &= x_1 - \frac{f(x_1)}{f^{'}(x_1)} \\
   &= 0.06242 - \frac{(0.06242)^3 - 0.165(0.06242)^2 + 3.993 \times 10^{-4}}{3(0.06242)^2 - 0.33(0.06242)} \\
   &= 0.06238
\end{aligned}$

$\therefore |\in_a| = \left|\frac{0.06238 - 0.06242}{0.06238}\right| \times 100 = 0.0716\%$

\textit{Iteration-3}:\\[0.5em]
$\begin{aligned}
   x_{3} &= x_2 - \frac{f(x_2)}{f^{'}(x_2)} \\
   &= 0.06238 - \frac{(0.06238)^3 - 0.165(0.06238)^2 + 3.993 \times 10^{-4}}{3(0.06238)^2 - 0.33(0.06238)} \\
   &= \red{0.06238}
\end{aligned}$

$\therefore |\in_a| = \left|\frac{0.06238 - 0.06238}{0.06238}\right| \times 100 = 0\%$
\end{multicols}



\section*{Secant Method}
\textbf{\textit{Question:}}
Use the Secant method to find the root of the equation $x^3 - 0.165x^2 + 3.993 \times 10^{-4} = 0$.\\
Use 0.02 and 0.055 as the initial guesses.

\begin{multicols}{2}
\noindent
\raggedcolumns
\textbf{\textit{Solution:}}\\[0.5em]
Let $f(x) = x^3 - 0.165x^2 + 3.993 \times 10^{-4}$\\[0.5em]
and $x_{-1} = 0.02$, \ $x_0 = 0.05$

\textit{Iteration-1}:\\[0.5em]
$\begin{aligned}
   x_{1} &= x_0 - \frac{f(x_0)(x_0 - x_{-1})}{f^(x_0)-f(x_{-1})} \\
   &= 0.05 - \frac{f(0.05)(0.05 - 0.02)}{f(0.05) - f(0.02)} \\
   &= 0.06242
\end{aligned}$

\columnbreak
\textit{Iteration-2}:\\[0.5em]
$\begin{aligned}
   x_{2} &= x_1 - \frac{f(x_1)}{f^{'}(x_1)} \\
   &= 0.06242 - \frac{(0.06242)^3 - 0.165(0.06242)^2 + 3.993 \times 10^{-4}}{3(0.06242)^2 - 0.33(0.06242)} \\
   &= 0.06238
\end{aligned}$

$\therefore |\in_a| = \left|\frac{0.06238 - 0.06242}{0.06238}\right| \times 100 = 0.0716\%$

\textit{Iteration-3}:\\[0.5em]
$\begin{aligned}
   x_{3} &= x_2 - \frac{f(x_2)}{f^{'}(x_2)} \\
   &= 0.06238 - \frac{(0.06238)^3 - 0.165(0.06238)^2 + 3.993 \times 10^{-4}}{3(0.06238)^2 - 0.33(0.06238)} \\
   &= \red{0.06238}
\end{aligned}$

$\therefore |\in_a| = \left|\frac{0.06238 - 0.06238}{0.06238}\right| \times 100 = 0\%$
\end{multicols}


\pagebreak

\section*{Trapezoidal Rule}
$h = \frac{b-a}{n} = \frac{30-8}{4} = 5.5$


$\begin{aligned}
   &\frac{b-a}{2n}\left[f(a) + 2\left\{\sum_{i=1}^{n-1}f(a+ih)\right\} + f(b)\right]
   \\
   = \ &\frac{30-8}{2\times 4}\left[f(a) + 2\left\{\sum_{i=1}^{4-1}f(a+ih)\right\} + f(b)\right]
   \\
   = \ &\frac{30-8}{2\times 4}\left[f(a) + 2f(a+h) + 2f(a+2h) + 2f(a+3h) + f(b)\right]
   \\
   = \ &\frac{30-8}{2 \times 4}\left[f(8) + 2f(8+5.5) + 2f(8+11) + 2f(8+16.5) + f(30)\right]
\end{aligned}$